\section{Pregunta de investigación e hipótesis}

\textbf{Pregunta de investigación.} ¿Puede una estrategia distribuida basada en consenso con un mecanismo adaptativo o robusto ofrecer un desempeño más consistente que una estrategia nominal en el transporte cooperativo multi-AGV cuando la carga presenta incertidumbre inercial y la comunicación entre agentes es únicamente local?

\textbf{Hipótesis de trabajo.}
\begin{enumerate}
  \item En escenarios con incertidumbre inercial de la carga, el controlador adaptativo o robusto mostrará menor error de seguimiento y menor error de formación que la línea base nominal en la mayor parte de los casos evaluados.
  \item La mejora anterior no será gratuita: se espera un posible incremento moderado del esfuerzo de control, pero sin pérdida global de estabilidad ni deterioro severo del tiempo de convergencia.
  \item La ventaja de la estrategia adaptativa o robusta será más visible cuando existan variaciones de masa, desplazamientos del centro de masas o perturbaciones externas acotadas, siempre que el grafo de comunicación conserve conectividad suficiente.
\end{enumerate}

\textbf{Criterio de contraste.} Las hipótesis se contrastarán mediante campañas de simulación reproducibles comparando ambas estrategias bajo el mismo conjunto de escenarios, condiciones iniciales y métricas de evaluación.
